\section{Resource Rational Navigation}

This section investigates when mice should rely on local heuristics vs global cognitive maps for navigation. Local navigation heuristics are rules of thumb (eg., the wall touching heuristic) that can be used to navigate without a detailed map of the environment. They are metabolically cheap and can be effective in simple environments, but they can lead to suboptimal behavior in more complex environments. In contrast, global cognitive maps are mental representations of the environment that allow for flexible navigation and planning. They are metabolically expensive to maintain and use, but they can lead to more optimal behavior in complex environments.
This extension focuses on finding the metabloic trade off between these two navigation strategies. We will use a resource rational analysis to determine when mice should rely on local heuristics vs global cognitive maps for navigation. This analysis will take into account the metabolic costs of both strategies, as well as the benefits of each strategy in different environments.
